\documentclass[journal]{IEEEtran}
\usepackage{cite}
\usepackage{amsmath,amssymb,amsfonts}
\usepackage{algorithmic}
\usepackage{graphicx}
\usepackage{textcomp}
\usepackage{xcolor}
\usepackage{booktabs}
\usepackage{multirow}
\usepackage{hyperref}

\begin{document}

\title{Climate-Aware Agricultural Commodity Price Forecasting: A Stacking Ensemble Approach with Uncertainty Quantification}

\author{Author 1, Author 2, Author 3%
\thanks{This work was supported by [Grant/Institution].}%
\thanks{The authors are with [Institution Name]. (e-mail: author@institution.edu).}}

\maketitle

\begin{abstract}
Agricultural commodity prices are highly susceptible to climate variability, yet traditional forecasting models often struggle to capture the complex, non-linear relationships between extreme weather events and market dynamics. In this paper, we propose a novel hybrid stacking ensemble approach that decomposes price prediction into temporal trend forecasting (via ARIMA) and non-linear climate-interaction modeling (via XGBoost). Using a comprehensive dataset from Tamil Nadu, India (2015-2024), we conduct per-commodity modeling across seven key staples and vegetables, avoiding the feature dominance pitfalls of global models. Furthermore, we apply split conformal prediction to generate distribution-free prediction intervals, providing risk-aware uncertainty quantification crucial for agricultural policy-making. Finally, we employ Granger causality to empirically validate the lagged temporal effects of rainfall deviation on price volatility. Our results demonstrate that the hybrid stacking ensemble significantly outperforms standalone baselines, particularly for volatile vegetables like onions and tomatoes, offering a robust framework for climate-resilient economic forecasting.
\end{abstract}

\begin{IEEEkeywords}
Machine Learning, Agricultural Economics, Time Series Forecasting, Conformal Prediction, Climate Change, Stacking Ensemble
\end{IEEEkeywords}

\section{Introduction}
\label{sec:introduction}
Agricultural markets are inherently volatile, driven by complex interactions between supply, demand, and increasingly, climate change. Extreme weather events such as droughts and excess monsoons directly disrupt crop yields, leading to severe price fluctuations. Accurate price forecasting is therefore critical for farmers to optimize harvesting strategies and for policymakers to ensure food security.

Traditional time-series models like Autoregressive Integrated Moving Average (ARIMA) effectively capture linear trends and seasonality but fail to account for exogenous shocks and non-linear climate interactions. Conversely, modern machine learning algorithms such as Random Forests and XGBoost excel at non-linear mappings but often struggle with long-term extrapolation of non-stationary trends. 

In this paper, we address these limitations through three primary contributions:
\begin{enumerate}
    \item \textbf{Hybrid Stacking Ensemble:} We introduce a two-level stacking architecture. Level 1 employs per-commodity ARIMA models to capture linear temporal dynamics. Level 2 utilizes an XGBoost meta-learner to integrate the ARIMA forecasts with lagged price features and climate variables (e.g., rainfall deviation).
    \item \textbf{Uncertainty Quantification:} We apply split conformal prediction to generate statistically guaranteed, distribution-free prediction intervals. This shifts the paradigm from simple point-forecasting to risk-aware decision support.
    \item \textbf{Empirical Causal Analysis:} We employ Granger causality testing to systematically validate the lagged temporal impact of extreme rainfall on specific commodity categories.
\end{enumerate}

\section{Methodology}
\label{sec:methodology}

\subsection{Data Integration and Preprocessing}
Our study integrates highly granular, multi-modal datasets spanning 2015 to 2024 for the state of Tamil Nadu, India. The primary target variable comprises daily retail prices for seven critical commodities: Rice, Wheat, Onion, Tomato, Potato, Sugar, and Milk. The exogenous features include district-level daily rainfall and its historical deviation. We aggregated the data to a monthly frequency to smooth extreme daily noise while preserving macroeconomic trends.

\subsection{The Commodity Dominance Problem}
A significant challenge in agricultural forecasting is "commodity dominance," where global models rely almost entirely on a categorical commodity identifier (often $>90\%$ feature importance), suppressing the learning of subtle climate signals. To resolve this, we strictly enforce \textit{per-commodity modeling}, ensuring the algorithms are forced to learn temporal and climate-driven dynamics rather than relying on baseline price disparities between different goods.

\subsection{Hybrid Stacking Architecture}
Let $y_{c,t}$ be the price of commodity $c$ at time $t$, and $X_{c,t}$ be the corresponding vector of climate features. 
In Level 1, we fit an ARIMA($p,d,q$) model for each commodity to generate a trend forecast $\hat{y}_{c,t}^{ARIMA}$.
In Level 2, an XGBoost meta-learner is trained to map the augmented feature space $Z_{c,t} = [X_{c,t}, y_{c,t-1}, y_{c,t-2}, y_{c,t-3}, \hat{y}_{c,t}^{ARIMA}]$ to the target $y_{c,t}$.

\subsection{Conformal Prediction}
To quantify uncertainty, we utilize split conformal prediction. The dataset is partitioned chronologically into training ($D_{train}$), calibration ($D_{cal}$), and testing ($D_{test}$) sets. We compute non-conformity scores $s_i = |y_i - \hat{y}_i|$ for all $i \in D_{cal}$. For a desired confidence level $1-\alpha$, we calculate the $\lceil (n_{cal}+1)(1-\alpha) \rceil / n_{cal}$ empirical quantile of $s$, denoted as $\hat{q}$. The prediction interval for a new test point is simply $[\hat{y}_{test} - \hat{q}, \hat{y}_{test} + \hat{q}]$.

\section{Results and Discussion}
\label{sec:results}

\subsection{Predictive Performance}
The models were evaluated using an expanding-window cross-validation strategy to simulate real-world forecasting constraints. The hybrid stacking ensemble demonstrated superior performance, effectively bridging the gap between ARIMA's trend-following and XGBoost's shock-responsiveness. 

% INSERT TABLE: tab:hybrid from latex_tables.tex

\subsection{Causal Analysis of Climate Impact}
Granger causality tests revealed distinct vulnerabilities across commodity types. Highly perishable vegetables (Tomato, Onion, Potato) exhibited statistically significant ($p < 0.05$) lagged causal relationships with rainfall deviation. Conversely, staples (Rice, Wheat) showed no significant Granger causality from climate variables, likely due to massive storage infrastructures and government interventions like Minimum Support Prices (MSP) that buffer climate shocks.

% INSERT TABLE: tab:granger from latex_tables.tex

\subsection{Uncertainty and Risk Assessment}
The conformal prediction intervals provided critical insights into market volatility. The empirical coverage rates aligned closely with the target confidence levels ($90\%$ and $95\%$), validating the distribution-free approach. The interval widths quantified inherent predictability: Sugar and Milk exhibited tight intervals (low volatility), whereas Tomatoes required significantly wider intervals to maintain coverage, reflecting their acute sensitivity to monsoon disruptions.

% INSERT TABLE: tab:conformal from latex_tables.tex

\section{Conclusion}
\label{sec:conclusion}
This study presents a robust, climate-aware forecasting framework for agricultural commodities. By decomposing the forecasting problem via a hybrid stacking ensemble and securing predictions with conformal intervals, we provide a holistic tool for agricultural risk management. Future work will integrate spatial graph neural networks to model cross-district supply chain contagions.

\bibliographystyle{IEEEtran}
\bibliography{references}

\end{document}
